\chapter{Pose3, Cant, and Spatial Railway State}

\section{Motivation}

\begin{remark}
Classical railway alignment models are based on planar geometry.
They describe the track centerline as a curve
\[
\gamma : s \mapsto (x(s),y(s)).
\]
\end{remark}

\begin{remark}
However, real railway systems are inherently three-dimensional.
In particular, the rotation of the track around its longitudinal axis
(cant) plays a fundamental role for vehicle dynamics.
\end{remark}

\begin{remark}
This motivates the introduction of an extended state description,
referred to as \emph{pose3}, which augments planar geometry by
rotational information.
\end{remark}

---

\section{Pose2 Revisited}

\begin{definition}[Pose2]
The planar pose of a curve is defined as
\[
\mathrm{pose2}(s)
=
\bigl(
\gamma(s), \theta(s)
\bigr),
\]
where
\[
\gamma(s)\in\R^2, \qquad \theta(s)\in\R.
\]
\end{definition}

\begin{remark}
The evolution of pose2 is governed by
\[
\theta'(s)=\kappa(s),
\qquad
\gamma'(s)=(\cos\theta(s),\sin\theta(s)).
\]
\end{remark}

---

\section{Cant as Rotational State}

\begin{definition}[Cant angle]
Let
\[
\phi : I \to \R
\]
denote the cant angle, describing rotation of the track cross-section
around the tangent direction.
\end{definition}

\begin{remark}
The cant angle represents the inclination of the track and directly
influences lateral acceleration and load distribution.
\end{remark}

\begin{definition}[Cant rate]
The derivative
\[
\phi'(s)
\]
describes the rate of change of cant along the track.
\end{definition}

---

\section{Pose3 Definition}

\begin{definition}[Pose3]
The spatial railway state (pose3) is defined as
\[
\mathrm{pose3}(s)
=
\bigl(
\gamma(s),
\theta(s),
\phi(s)
\bigr).
\]
\end{definition}

\begin{remark}
Pose3 extends pose2 by adding a rotational degree of freedom.
It captures both geometric and dynamic-relevant aspects of the track.
\end{remark}

---

\section{State Evolution}

\begin{proposition}[Pose3 evolution]
The pose3 state satisfies
\[
\begin{aligned}
\gamma'(s) &= (\cos\theta(s),\sin\theta(s)), \\
\theta'(s) &= \kappa(s), \\
\phi'(s)   &= \omega(s),
\end{aligned}
\]
where \(\omega(s)\) denotes the cant-rate control.
\end{proposition}

\begin{remark}
Thus, pose3 is governed by two control functions:
\[
\kappa(s) \quad \text{and} \quad \omega(s).
\]
\end{remark}

---

\section{Reconstruction Principle}

\begin{proposition}
Given
\[
\kappa(s), \quad \omega(s),
\]
and initial conditions
\[
\gamma(0), \quad \theta(0), \quad \phi(0),
\]
the pose3 state is uniquely determined.
\end{proposition}

\begin{proof}
The result follows by integration of the differential system.
\end{proof}

---

\section{Relation to Physical Quantities}

\subsection{Effective Lateral Acceleration}

\begin{definition}
The effective lateral acceleration is
\[
a_{\mathrm{eff}}(s)
=
v^2 \kappa(s) - g \sin\phi(s).
\]
\end{definition}

\begin{remark}
This quantity determines passenger comfort and safety.
\end{remark}

---

\subsection{Jerk}

\begin{definition}
The jerk is given by
\[
j(s)
=
v^3 \kappa'(s) - g v \cos\phi(s)\,\phi'(s).
\]
\end{definition}

\begin{remark}
The jerk measures the rate of change of acceleration and is a key
criterion for comfort.
\end{remark}

---

\section{Interpretation as Control System}

\begin{remark}
The pose3 model can be interpreted as a control system:
\[
\begin{aligned}
\text{state:} &\quad (\gamma,\theta,\phi), \\
\text{controls:} &\quad \kappa, \ \omega.
\end{aligned}
\]
\end{remark}

\begin{remark}
This formulation provides a natural bridge to optimal control and
optimization methods.
\end{remark}

---

\section{Cant and Curvature Coupling}

\begin{remark}
Curvature and cant are not independent.
They jointly determine dynamic behaviour.
\end{remark}

\begin{definition}[Equilibrium condition]
A perfectly balanced design satisfies
\[
v^2 \kappa(s) = g \sin\phi(s).
\]
\end{definition}

\begin{remark}
In practice, deviations from this condition are allowed but must remain
within regulatory limits.
\end{remark}

---

\section{Towards Schwerpunkt-Trassierung}

\begin{remark}
The pose3 formulation enables a reinterpretation of alignment design.
Instead of designing geometry alone, one designs a spatial state
trajectory.
\end{remark}

\begin{remark}
In particular, the concept of \emph{Schwerpunkt-Trassierung} can be
understood as the design of a trajectory that optimally balances
curvature and cant in a dynamic sense.
\end{remark}

\RESEARCH{Formal definition of Schwerpunkt-based alignment design.}

\OPEN{Integration of vehicle dynamics into the pose3 framework.}

---

\section{Relation to Existing Standards}

\begin{remark}
Existing railway design standards typically prescribe:
\begin{itemize}
\item limits on curvature,
\item limits on cant,
\item limits on cant gradient.
\end{itemize}
\end{remark}

\begin{remark}
These rules can be interpreted as constraints on the pose3 state and
its derivatives.
\end{remark}

---

\section{Connection to the AIM Framework}

\begin{summary}
The introduction of pose3 extends alignment modelling from planar
geometry to a spatial state description.

This provides:
\begin{itemize}
\item a unified representation of geometry and dynamics,
\item a natural interface to optimization methods,
\item a foundation for advanced concepts such as
Schwerpunkt-Trassierung.
\end{itemize}

Within the AIM framework, pose3 serves as the central state model for
railway alignment design.
\end{summary}
